\documentclass[10pt,a4paper]{amsart}
\usepackage[margin=19mm]{geometry}
\usepackage{amsmath,amssymb,mathtools}
\usepackage[scr=rsfs,cal=euler]{mathalfa}
\usepackage{xfrac}
\usepackage[unicode,hidelinks,bookmarks=false]{hyperref}
\newcommand{\ie}{i.e.\ }
\newcommand{\cf}{cf.\ }
\newcommand{\resp}{respectively\ }
\newcommand{\Subsection}{\subsection}
\providecommand{\nobreakdash}{\nobreak\hbox{-}}
\setlength{\emergencystretch}{2em}
\multlinegap0pt
\usepackage{bm}
\usepackage{bbm}
\usepackage{dsfont}
\usepackage{xspace}
\usepackage{cancel}
\usepackage{mathtools}
\usepackage{amsthm}
\makeatletter
\@ifundefined{thm}{\newtheorem{thm}{Theorem}}{}
\@ifundefined{lem}{\newtheorem{lem}[thm]{Lemma}}{}
\makeatother
\title[On the $m$-graph of a finite Abelian Group]{On the $m$-graph of a finite Abelian Group}
\author{Ayman Badawi}
\date{Source: arXiv:2601.10830v1 (CC BY 4.0)}
\begin{document}
\maketitle

\noindent\emph{Source and licence.} On the $m$-graph of a finite Abelian Group. Version arXiv:2601.10830v1 (CC BY 4.0), distributed under the Creative Commons Attribution 4.0 International licence, as recorded in the arXiv licence metadata for this exact version and shown on the abstract page https://arxiv.org/abs/2601.10830v1. Full rights notice: licenses/arxiv-2601.10830-CC-BY-4.0.txt.\par\smallskip

\noindent\textit{Editorial scope.} Counted unit: the Thm whose source label is \texttt{T5} in the article (source0.tex lines 311--322, source label \texttt{T5}), delivered together with the article's own complete original proof of it (source0.tex lines 324--346, one \texttt{proof} environment of 23 lines). No mathematics is written, repaired or re-derived: statement and proof are the article's own text line for line. Required context retained uncounted: T1 (lines 45--51); L1 (lines 86--88). Every definition, display and label of those lines is retained unchanged. Omitted from this package and not redistributed: 1--44, 52--85, 89--310, 323--323, 347--458. Nothing inside a retained excerpt is omitted or altered: every retained body is delivered exactly as the article's lines are written, so no omission receipt of any kind accompanies this package. Citations: the retained excerpts carry no citation command at all, so no bibliography entry of the article is transcribed and no reference is redistributed. Reference adaptation: none was needed and none was made; every reference inside the delivered unit resolves inside it, so no unresolved reference or citation is produced. Printed slips kept literal: no printed word, symbol, hypothesis or number of the counted statement or of its proof was altered. Layout and class: the article's own class is \texttt{amsart} with packages including amssymb, amsfonts, graphicx; the wrapper is the lane's pinned \texttt{amsart} editorial wrapper at 10pt on A4, which restates the theorem-like environments the retained text needs and adds the article's own macro definitions that the retained text needs (none), with extra wrapper packages (bm, bbm, dsfont, xspace, cancel, mathtools). The delivered numbering is the unit's own. Assets: no retained span contains a figure environment, a graphics-inclusion command, a PDF-page inclusion, a table or a TikZ picture; no image file is copied into this package, the package declares no asset dependency and no image file is redistributed with it. Licence: version arXiv:2601.10830v1 of this article is distributed under the Creative Commons Attribution 4.0 International licence, as recorded in the arXiv licence metadata for this exact version and shown on the abstract page https://arxiv.org/abs/2601.10830v1. A full rights notice is delivered at licenses/arxiv-2601.10830-CC-BY-4.0.txt. Characters: every retained excerpt is pure ASCII, so the delivered engine, XeLaTeX with its default Latin Modern text font, reports no missing character.
\begin{thm}\label{T1}
	Let $H$ be a finite abelian group of order $n \geq 2$, $m > 1$ be an integer, and $k = gcd(m, n)$. The following statements are equivalent.
	\begin{enumerate}
		\item $m-G(H)$ is connected.
		\item $k\mid n$ and if $p$ is a prime factor of $n$, then $p \mid k$ (and hence $p \mid m$).
	\end{enumerate}
\end{thm}

\begin{lem}\label{L1}
	Let $n, m > 1$ be integers and assume that the equation $mx = a$ has a solution in $Z_n$ for some $a\in Z_n$. Let $k = gcd(m, n) $. Then $k \mid a$ in $Z$ and the equation $mx = a$ in $Z_n$ has exactly $k$ distinct solutions in $Z_n$.  Furthermore, assume that every prime factor of $n$ is a prime factor of $m$ and $a \not = 0$. Then $ma \not = a$ in $Z_n$.	
	\end{lem}

\begin{thm}\label{T5}
	Let $n_1, ..., n_i$  be positive integers such that $i \geq 2$,  $gcd(n_e, n_j) \not = 1$ for every $e \not = j$, $1 \leq e, j \leq i$, and $n = n_1n_2 \cdots n_i$. Let $m$ be a positive integer such that every prime factor of $n$ is a prime factor of $m$, and $k = gcd(n, m)$. For each $1 \leq j \leq i$, let $d_j = gcd(m, n_j)$ and $n_j = q_jkd_j$ for some positive integer $q_j$. Let $H$ be the abelian group $Z_{n_1}  \times \cdots \times Z_{n_i}$. Then the following statements hold.
	\begin{enumerate}
		\item  $H$ is not cyclic and the $m$-$G(H)$ is connected.
		\item $deg((0, ..., 0)) = d_1d_2\cdots d_i  - 1 \geq 2^i -1$.
		\item Let $(a_1, ..., a_i) \in H$. If $d_j \nmid a_j$ (in $Z$) for some $1 \leq j \leq i$, then $deg((a_1, ..., a_i)) = 1$.
		\item Let $(a_1, ..., a_i) \in H \setminus \{ (0, ..., 0)\}$. If $d_j \mid a_j$ for every $1 \leq j \leq i$, then $deg((a_1, ... , a_i) = d_1d_2 \cdots d_i + 1$.
		\item There are exactly $q_1q_2 \cdots q_i - 1$ distinct elements in $H$ of degree $d_1d_2 \cdots d_i + 1$.
		\item There are exactly $n - q_1q_2 \cdots q_i$ distinct elements in $H$ of degree 1.
		\item The $m$-$G(H)$ is a tree.
	\end{enumerate}
\end{thm}

\begin{proof}
	 First, for every $1 \leq j \leq i$, observe that $d_j \geq 2$ and every prime factor of $n_j$ is a prime factor of $d_j$.
	 \vskip0.1in
	{\bf (1)}. It is clear that $k \mid n$ and every prime factor of $n$ is a prime factor of $k$. Hence by Theorem \ref{T1}, the $m$-$G(H)$ is connected. Since the $gcd(n_e, n_j) \not = 1$ for every $e \not = j$, $1 \leq e, j \leq i$, we conclude that $H$ is not cyclic.
	\vskip0.1in
	{\bf (2)} The equation $mX = m(x_1,..., x_i) = (mx_1, ... , mx_i) = (0, ..., 0)$ has a solution in $H$ if and only if $(d_1x_1, ... , d_ix_i) = (0, ... , 0)$ has a solution in $H$ by Lemma \ref{L1}. For each $1 \leq j \leq i$, the equation $d_jx_j = 0$ has exactly $d_j \geq 2$ distinct solutions (including 0) in $Z_{n_j}$ by Lemma \ref{L1}. Thus the number of distinct nonzero elements that are adjacent to $(0, ... , 0)$ is $d_1d_2 \cdots d_i - 1$. Since $d_j \geq 2$ for every $1\leq j \leq i$ and $i \geq 2$, we have $deg((0, ... , 0)) = d_1d_2 \cdots d_i - 1 \geq 2^i - 1$.

{\bf (3)}. Let $(a_1, ..., a_i) \in H$ and assume that $d_j \nmid a_j$ (in $Z$) for some $1 \leq j \leq i$. Hence $(a_1, ... , a_i) \not = (0, ... , 0)$. Since $d_jx_j = a_j$ has no solution in $Z_{n_j}$ by Lemma \ref{L1}, we conclude that $mX = (mx_1, ..., mdx_i) = (a_1, ... , a_i)$ has no solution in $H$. Since $d_j \nmid a_j$, $a_j \not = 0$. Thus $m(a_1, ..., a_j, ..., a_i) \not = (a_1, ..., a_j, ... , a_i)$ by Lemma \ref{L1}. Hence $m(a_1, ..., a_j, ..., a_i)$ is the only element in $H$ that is adjacent to $(a_1, ..., a_j, ..., a_i)$. Thus $deg((a_1, ..., a_j, ..., a_i)) = 1$.
\vskip0.1in
{\bf (4)}. Let $(a_1, ..., a_i) \in H \setminus \{ (0, ..., 0)\}$ and assume that $d_j \mid a_j$ for every $1 \leq j \leq i$. The equation $mX = m(x_1,..., x_i) = (mx_1, ... , mx_i) = (a_1, ..., a_i)$ has a solution in $H$ if and only if $(d_1x_1, ... , d_ix_i) = (a_1, ... , a_i)$ has a solution in $H$ by Lemma \ref{L1}. For each $1 \leq j \leq i$, the equation $d_jx_j = a_j$ has exactly $d_j$ distinct solutions  in $Z_{n_j}$ by Lemma \ref{L1}. Thus there are $d_1d_2 \cdots d_i$ distinct elements in $H$ that are adjacent to $(a_1, ... , a_i)$. Since $m(a_1, ..., a_i) \not = (a_1, ..., a_i)$ by Lemma \ref{L1}, we conclude that $m(a_1, ..., a_i)$ is adjacent to $(a_1, ..., a_i)$. Thus $deg((a_1, ... , a_i) = d_1d_2 \cdots d_i + 1$.
\vskip 0.1in
{\bf (5)}. Let $L = \{(a_1, ..., a_i) \in H \mid \   d_j \mid a_j$ \ for every $1 \leq j \leq i\}$. Let $w \in L$, then $w = (c_1d_1, ..., c_id_i)$ such that for each $1 \leq j \leq i$, we have $1 \leq c_j \leq q_j$. Hence $|L| = q_1q_2 \cdots q_i$. Thus $|L \setminus (0, ... , 0)| = q_1q_2 \cdots q_i - 1$. Thus there are exactly $q_1q_2 \cdots q_i - 1$ distinct elements in $H$ of degree $d_1d_2 \cdots d_i + 1$.

\vskip 0.1in
{\bf (6)}.  Let $L = \{(a_1, ..., a_i) \in H \mid \   d_j \mid a_j$ \ for every $1 \leq j \leq i\}$ and $U = \{(a_1, ..., a_i) \in H \mid \   d_j \nmid a_j$ \ for some $1 \leq j \leq i\}$. It is clear that $L \cup U = H$ and $L \cap U = \{\}$. Since $|L| = q_1q_2 \cdots q_i$ by (5) and $|H| = n$, we conclude that $|U| = n - |L| = n - q_1q_2 \cdots q_i$.  Since each element in $U$ is of degree one by (3), we conclude that there are exactly $n - q_1q_2 \cdots q_i$ distinct elements in $H$ of degree 1.

\vskip 0.1in
{\bf (7)}. Let $E$ be the set of all edges of $m$-$G(H)$. We show $|E| = |H| - 1 = n - 1$. It is known that $|E| = \frac{\sum_{v \in H} deg(v)}{2}$. Let $L$ and $U$ as in (6). Then by (2) - (6), we have
$$\sum_{v \in H} deg(v) = \sum_{v \in L} deg(v) + \sum_{v \in U} deg(v) = (q_1q_2 \cdots q_i - 1)(d_1d_2 \cdots d_i + 1) + 1(d_1d_2\cdots d_i - 1) + $$\\ $$(n - q_1q_2 \cdots q_i)1 = q_1q_2 \cdots q_i d_1d_2\cdots d_i + q_1q_2 \cdots q_i - d_1d_2\cdots d_i - 1 + d_1d_2\cdots d_i - 1 +$$ \newline   $$n - q_1q_2\cdots q_i = 2n - 2$$.

Hence $|E| = \frac{2n - 2}{2} = |H| -1 = n-1$. Since the $m$-$G(H)$ is connected and $|E| = |H| - 1|$, we conclude that the $m$-$G(H)$ is a tree.

\end{proof}

\end{document}
